% Class-compatibility preamble retained for a reader excerpt of
%   Lisa Sauermann, Dmitrii Zakharov, "A Sharp Ramsey Theorem for Ordered
%   Hypergraph Matchings", arXiv:2309.04813v4, CC BY 4.0.
%
% The source pins the publisher class aic (the AIMS journal class), which is not part of
% the delivery baseline, so this file supplies the amsart baseline and then copies the
% author's own declaration block mechanically and unchanged
% (pinned-original-source.tex lines 25 and 28-59: the ams packages, the author's eleven
% newtheorem declarations with [theorem] as the shared counter, the definition theorem
% style, the author's comment that the file should add as few macro or package
% definitions as possible, and the author's own macros \su, \sm, \Z, \HH, \conv and \PP).
% The author's tikz package (pinned source line 26) serves the three figures of the paper
% that this excerpt does not carry and is therefore not reproduced; the author's babel,
% fontenc and inputenc lines (pinned source lines 47-49) are layout/font choices of the
% publisher build and are not part of the mathematical excerpt, the delivery engine being
% XeLaTeX with UTF-8 input.
% No mathematical macro, statement, symbol or assumption is introduced, changed or
% removed by this file.
\documentclass[10pt,a4paper]{amsart}
\usepackage{amsthm,amsfonts,amssymb,amsmath}
\usepackage{graphicx}

\newtheorem{theorem}{Theorem}[section]
\newtheorem{proposition}[theorem]{Proposition}
\newtheorem{lemma}[theorem]{Lemma}
\newtheorem{conjecture}[theorem]{Conjecture}
\newtheorem{fact}[theorem]{Fact}
\newtheorem{claim}[theorem]{Claim}
\newtheorem{corollary}[theorem]{Corollary}
\newtheorem{definition}[theorem]{Definition}
\newtheorem{problem}[theorem]{Problem}
\newtheorem{question}[theorem]{Question}
\newtheorem{observation}[theorem]{Observation}

\theoremstyle{definition}
\newtheorem{remark}[theorem]{Remark}

% The author notes at this point that as few macro or package
% definitions as possible should be added; the six macros below are the
% author's own and are copied unchanged.

\newcommand{\su}{\subseteq}
\newcommand{\sm}{\setminus}
\newcommand{\Z}{\mathbb{Z}}
\newcommand{\HH}{\mathcal{H}}
\newcommand{\conv}{\operatorname{conv}}
\newcommand{\PP}{\mathcal{P}}
\usepackage{cleveref}
\crefname{originalsection}{original section}{original sections}
\Crefname{originalsection}{Original section}{Original sections}
\makeatletter
\newlabel{sect-preliminaries}{{2}{source}{}{sourcepointer-1}{}}
\newlabel{sect-preliminaries@cref}{{[originalsection][0][]2}{[1][0][]source}{}{}{}}
\makeatother
\crefname{originaltheorem}{original theorem}{original theorems}
\Crefname{originaltheorem}{Original theorem}{Original theorems}
\makeatletter
\newlabel{thm_es}{{2.1}{source}{}{sourcepointer-2}{}}
\newlabel{thm_es@cref}{{[originaltheorem][0][]2.1}{[1][0][]source}{}{}{}}
\makeatother
\crefname{originaltheorem}{original theorem}{original theorems}
\Crefname{originaltheorem}{Original theorem}{Original theorems}
\makeatletter
\newlabel{thm-paramteters-upper-bound}{{1.2}{source}{}{sourcepointer-3}{}}
\newlabel{thm-paramteters-upper-bound@cref}{{[originaltheorem][0][]1.2}{[1][0][]source}{}{}{}}
\makeatother
\crefname{originaltheorem}{original theorem}{original theorems}
\Crefname{originaltheorem}{Original theorem}{Original theorems}
\makeatletter
\newlabel{thm-paramteters-lower-bound-stronger}{{4.3}{source}{}{sourcepointer-4}{}}
\newlabel{thm-paramteters-lower-bound-stronger@cref}{{[originaltheorem][0][]4.3}{[1][0][]source}{}{}{}}
\makeatother
\crefname{originalequation}{original equation}{original equations}
\Crefname{originalequation}{Original equation}{Original equations}
\makeatletter
\newlabel{Lp}{{4.1}{source}{}{sourcepointer-5}{}}
\newlabel{Lp@cref}{{[originalequation][0][]4.1}{[1][0][]source}{}{}{}}
\makeatother
\title{An r-partite ordered hypergraph matching realising prescribed clique-size bounds via a monotone point set}
\author{Lisa Sauermann, Dmitrii Zakharov}
\date{arxiv:2309.04813v4; CC BY 4.0}
\begin{document}
\maketitle
\noindent\textit{Editorial note.} All delivered mathematical text is the author's own, copied from the pinned author source sources/combinatorics-fl02-a-sharp-ramsey-theorem-for-ordered/pinned-original-source.tex; no statement or proof step was rewritten, reordered or removed.\par\medskip
\noindent\textit{Editorial note.} Preamble provenance: the source pins the publisher class aic, which is not part of the delivery baseline, so this card uses sources/combinatorics-fl02-a-sharp-ramsey-theorem-for-ordered/retained-preamble.tex, which supplies the amsart baseline and then copies the author's ams packages, the author's eleven newtheorem declarations sharing the theorem counter, the definition theorem style and the author's own macros \textbackslash{}su, \textbackslash{}sm, \textbackslash{}Z, \textbackslash{}HH, \textbackslash{}conv and \textbackslash{}PP (pinned source lines 25 and 28-59) verbatim. No macro was invented and no mathematical symbol was changed.\par\medskip
\noindent\textit{Editorial note.} Bibliography provenance: the citations of the counted proof are resolved with sources/combinatorics-fl02-a-sharp-ramsey-theorem-for-ordered/bibliography.tex, an excerpt of the author's own thebibliography environment (pinned source lines 598-703) copied character for character with all eighteen entries kept; no entry was rewritten, reordered, merged, split or invented and no citation key was changed.\par\medskip
\noindent\textit{Editorial note.} External references are navigation pointers only; every deliberately kept step of the author's argument stays inside the counted spans. Their printed numbers are the author's own numbering under the author's newtheorem declarations, computed mechanically by the lane's scanner.\par\medskip
\noindent\textit{Editorial note.} The labels the counted proof invokes outside its own span, sect-preliminaries, thm\_es, thm-paramteters-upper-bound and thm-paramteters-lower-bound-stronger, are navigation-pointer stubs; thm-paramteters-lower-bound-stronger has its own card in this lane.\par\medskip
\noindent\textit{Allowed standard prerequisites (editorial).}\begin{itemize}
\item The multi-dimensional Erdos-Szekeres theorem in its second form, producing a point set with bounded monotone subsequences. The author's theorem labelled thm\_es (source lines 232-238), which the counted proof invokes by label. Produces the point collection Z whose monotone sequences have the prescribed bounded lengths.
\item The correspondence between r-patterns of the finest ordered partition and monotone functions, and the quantity L\_P of a collection. Established at the end of the author's preliminaries (discussion at the end of the section, source lines 281-288) and used in the counted proof. Turns the bounded monotone sequences of the point set into the clique-size bounds L\_P of the constructed matching.
\end{itemize}\medskip
\ifdefined\hypertarget\hypertarget{sourcepointer-1}{}\fi
\noindent\textit{Source reference.} Original section 2: sources/combinatorics-fl02-a-sharp-ramsey-theorem-for-ordered/pinned-original-source.tex, source lines 226--226. Navigation pointer only: it locates the author's preliminaries where the pattern parameterisation is set up; no text is reproduced.\par\medskip
\ifdefined\hypertarget\hypertarget{sourcepointer-2}{}\fi
\noindent\textit{Source reference.} Original theorem 2.1: sources/combinatorics-fl02-a-sharp-ramsey-theorem-for-ordered/pinned-original-source.tex, source lines 232--232. Navigation pointer only: the counted proof applies this multi-dimensional Erdos-Szekeres theorem; it is declared as a named standard prerequisite and is not reproduced.\par\medskip
\ifdefined\hypertarget\hypertarget{sourcepointer-3}{}\fi
\noindent\textit{Source reference.} Original theorem 1.2: sources/combinatorics-fl02-a-sharp-ramsey-theorem-for-ordered/pinned-original-source.tex, source lines 199--199. Navigation pointer only: the counted proof refers to the author's upper-bound theorem in the surrounding discussion; the theorem is not reproduced.\par\medskip
\ifdefined\hypertarget\hypertarget{sourcepointer-4}{}\fi
\noindent\textit{Source reference.} Original theorem 4.3: sources/combinatorics-fl02-a-sharp-ramsey-theorem-for-ordered/pinned-original-source.tex, source lines 455--455. Navigation pointer only: the counted construction is the building block of that theorem, which has its own card in this lane.\par\medskip
\ifdefined\hypertarget\hypertarget{sourcepointer-5}{}\fi
\noindent\textit{Source reference.} Original equation 4.1: sources/combinatorics-fl02-a-sharp-ramsey-theorem-for-ordered/pinned-original-source.tex, source lines 549--549. Navigation pointer only: it labels the author's display defining L\_P, used by the counted proof; the display is not reproduced.\par\medskip
\begin{lemma}\label{obs-construction-partite}
    Let $r$ be a positive integer, and let $\lambda^{(r)}=(1,\dots,1)$ be the unique ordered partition of $r$ into $r$ parts. For every $r$-pattern $P\in \PP(\lambda^{(r)})$, consider a positive integer $m_P$. Then there exists an $r$-partite ordered $r$-uniform hypergraph matching $\HH$ of size
\[|\HH|=\prod_{P\in \PP(\lambda^{(r)})} m_{P},\]
such that $L_P(\HH)\le m_P$ for all $r$-patterns $P\in \PP(\lambda^{(r)})$ and $L_P(\HH)=1$ for all other $r$-patterns $P\not\in \PP(\lambda^{(r)})$.
\end{lemma}

\begin{proof}[Proof of Lemma~\ref{obs-construction-partite}]
    Recall from the discussion at the end of Section~\ref{sect-preliminaries} that any $r$-pattern $P\in \PP(\lambda^{(r)})$ is of the form $P=P(\lambda^{(r)},\tau)$ for a unique function $\tau:\{1,\dots,r\}\to \{1,-1\}$ with $\tau(1)=1$. For every such function $\tau$, let us define $m_\tau=m_{P(\lambda^{(r)},\tau)}$. By the second part of Theorem~\ref{thm_es}, there exists a collection of points $Z\su \mathbb{R}^r$ of size $|Z|=\prod_{\tau} m_{\tau}=\prod_{P\in \PP(\lambda^{(r)})} m_{P}$ such that the $|Z|\cdot r$ coordinates of all the points in $Z$ are all distinct, and such that for every function $\tau:\{1,\dots,r\}\to \{1,-1\}$ with $\tau(1)=1$, every $\tau$-monotone sequence in $Z$ has length at most $m_\tau=m_{P(\lambda^{(r)},\tau)}$.
    
    Without loss of generality, we may assume that all points in $Z$ have positive coordinates (otherwise, we can translate the set $Z$ accordingly). Let $N>0$ be a number larger than all coordinates of all points in $Z$. Now, for every point $z=(z_1,\dots,z_r)\in Z\su \mathbb{R}^r$, let us define $e(z)=\{N+ z_1, 2N + z_2, \ldots, r N+z_r \}\su \mathbb{R}$, and note that the sets $e(z)$ for $z\in Z$ are pairwise disjoint. Taking $\HH=\{e(z)\mid z\in Z\}$ to be the collection of the sets $e(z)$ for all $z\in Z$, we obtain an ordered $r$-uniform hypergraph matching $\HH$ by considering the natural ordering on $\mathbb{R}$ on the vertex set. Note that $\HH$ has size $|\HH|=|Z|=\prod_{P\in \PP(\lambda^{(r)})} m_{P}$. Furthermore, $\HH$ is $r$-partite (by the choice of $N$) and so in particular we have $L_{P}(\HH)=1$ for any $P\not\in \PP(\lambda^{(r)})$. It remains to check $L_{P}(\HH)\le m_P$ for every $P\in \PP(\lambda^{(r)})$.

    So let $P\in \PP(\lambda^{(r)})$ and let $\tau:\{1,\dots,r\}\to \{1,-1\}$ with $\tau(1)=1$ be such that $P=P(\lambda^{(r)},\tau)$, then we have $m_\tau=m_P$. Let $z^{(1)},\dots,z^{(t)}\in Z$ be such that the edges $e(z^{(1)}),\dots,e(z^{(t)})$ form $P$-clique in $\HH$. Then for any $1\le \ell<\ell'\le t$, for any $i\in \{1,\dots,r\}$ with $\tau(i)=1$, the $i$-th vertex $iN+z^{(\ell)}_i$ of $e(z^{(\ell)})$ comes before the $i$-th vertex $iN+z^{(\ell')}_i$ of $e(z^{(\ell')})$ and so we have $z^{(\ell)}_i<z^{(\ell')}_i$. Similarly for any $i\in \{1,\dots,r\}$ with $\tau(i)=1$, the $i$-th vertex of $e(z^{(\ell)})$ comes after the $i$-th vertex of $e(z^{(\ell')})$ and so we have $z^{(\ell)}_i>z^{(\ell')}_i$. Thus, for every $i\in \{1,\dots,r\}$ with $\tau(i)=1$, the $i$-th coordinates of the points $z^{(1)},\dots,z^{(t)}$ form an increasing sequence 
    $z^{(1)}_1<\dots<z^{(t)}_i$, and for every $i\in \{1,\dots,r\}$ with $\tau(i)=-1$, they form a decreasing sequence $z^{(1)}_1>\dots>z^{(t)}_i$. Thus, $z^{(1)},\dots,z^{(t)}$ is a $\tau$-monotone sequence of points in $Z$ and hence $t\le m_\tau\le m_P$. This shows that $L_P(\HH)\le m_P$, as desired.
\end{proof}

% Bibliography excerpt retained for reader excerpts of
%   Lisa Sauermann, Dmitrii Zakharov, "A Sharp Ramsey Theorem for Ordered Hypergraph Matchings",
%   arXiv:2309.04813v4, CC BY 4.0.
%
% Every entry below is copied character for character from the author's own thebibliography
% environment of the pinned author source (lines 598-703).  No entry was rewritten, reordered,
% merged, split or invented and no citation key was changed; all 18 entries are kept so that
% any citation of a delivered excerpt resolves without guessing an entry.
\begin{thebibliography}{99}

%\bibitem{bergelson-johnson-moreira}
%Vitaly Bergelson, John H. Johnson Jr., and Joel Moreira.
%\newblock New polynomial and multidimensional extensions of classical partition
%  results.
%\newblock 2015, arXiv:1501.02408.

%\bibitem{cilleruelo}
%Javier Cilleruelo.
%\newblock Combinatorial problems in finite fields and {S}idon sets.
%\newblock {\em Combinatorica}, 32(5):497--511, 2012.


\bibitem{AJKS} 
Michael Anastos, Zhihan Jin, Matthew Kwan, and Benny Sudakov.
\newblock Extremal, enumerative and probabilistic results on ordered hypergraph matchings.
\newblock {\em Forum of Mathematics Sigma} 13: Paper No.~e55, 39~pp, 2025.

\bibitem{Campos} 
Marcelo Campos, Simon Griffiths, Robert Morris, and Julian Sahasrabudhe. 
\newblock An exponential improvement for diagonal Ramsey.
\newblock 2023, arXiv:2303.09521.

\bibitem{Caro} 
Yair Caro. 
\newblock New Results on the Independence Number.
\newblock Technical Report, Tel-Aviv University, 1979.

\bibitem{ExpES} 
Recep Altar Çiçeksiz, Recep Altar, Zhihan Jin, Eero Räty, and István Tomon. 
\newblock Exponential Erd\H{o}s-Szekeres theorem for matrices.
\newblock {\em Proceedings of the London Mathematical Society} 129, no. 3: Paper No. e12632, 14~pp, 2024.

\bibitem{Dil} 
Robert P. Dilworth 
\newblock A Decomposition Theorem for Partially Ordered Sets.
\newblock {\em Annals of Mathematics}  51: 161--166, 1950.

\bibitem{DGR22}
Andrzej Dudek, Jarosław Grytczuk, and Andrzej Ruci\'{n}ski. 
\newblock Ordered unavoidable sub-structures in matchings and random matchings.
\newblock {\em Electronic Journal of Combinatorics} 31, no. 2: Paper No. 2.15, 27 pp, 2024

\bibitem{DGR} 
Andrzej Dudek, Jarosław Grytczuk, and Andrzej Ruci\'{n}ski. 
\newblock Erd\H{o}s-Szekeres type Theorems for ordered uniform matchings.
\newblock {\em Journal of Combinatorial Theory, Series B} 170: 225--259, 2025.

\bibitem{ES} 
Paul Erd{\H o}s and George Szekeres. 
\newblock A combinatorial problem in geometry.
\newblock {\em Compositio Mathematica} 2: 463--470, 1935.

\bibitem{GL} Timothy W. Gowers and Jason Long.
\newblock The length of an s-increasing sequence of r-tuples. \newblock {\em Combinatorics, Probability and Computing} 30, no. 5: 686--721, 2021.

\bibitem{graham-rothschild-spencer} 
Ronald L. Graham, Bruce L. Rothschild, and Joel H. Spencer. 
\newblock {\em Ramsey theory.} Vol. 20. John Wiley \& Sons, 1991.


\bibitem{Lee} 
Choongbum Lee. 
\newblock Ramsey numbers of degenerate graphs.
\newblock {\em Annals of Mathematics} 185, no. 3: 791--829, 2017.

\bibitem{LS} 
Nathan Linial and Michael Simkin. 
\newblock Monotone subsequences in high-dimensional permutations. 
\newblock {\em Combinatorics, Probability and Computing} 27, no. 1: 69--83, 2018.

\bibitem{MV} 
Sam Mattheus and Jacques Verstraete. 
\newblock The asymptotics of $r(4,t)$.
\newblock {\em Annals of Mathematics} 199, no.~2: 919–941,  2024.

%\bibitem{Mirsky} 
%Leon Mirsky. 
%\newblock A dual of Dilworth’s decomposition theorem.
%\newblock {\em Amer. Math. Monthly} 78, no. 8: 876-877, 1971.

\bibitem{MS} 
Guy Moshkovitz and Asaf Shapira. 
\newblock Ramsey theory, integer partitions and a new proof of the Erdős–Szekeres theorem. 
\newblock {\em Advances in Mathematics} 262: 1107--1129, 2014.

\bibitem{promel} 
Hans Jürgen Prömel. 
\newblock {\em Ramsey theory for discrete structures.} Vol. 3. New York: Springer, 2013.

\bibitem{RR} 
Christian Reiher, and Vojt\v{e}ch R\"odl. 
\newblock The girth Ramsey theorem.
\newblock 2023, arXiv:2308.15589.

\bibitem{ST} 
Tibor Szab\'{o} and G\'{a}bor Tardos. 
\newblock A multidimensional generalization of the Erd\H{o}s–Szekeres lemma on monotone subsequences.
\newblock {\em Combinatorics, Probability and Computing} 10, no. 6: 557--565, 2001.

\bibitem{Wei} 
Victor K. Wei. 
\newblock A Lower Bound on the Stability Number of a Simple Graph.
\newblock {\em Technical memorandum}, TM 81-11217-9, Bell Laboratories, 1981.


\end{thebibliography}
\end{document}
